% الجزء: نظرية البرهان
% الفصل: حذف القطع
% القسم: ce-largest

\documentclass[../../../include/open-logic-section]{subfiles}

\begin{document}

\olfileid{pt}{cut}{inv}

\olsection{في إزالة القطوع العظمى رتبة}

يستفاد من برهان \olref[seq][inv]{lem:G3c-invert} الأمر الآتي: إذا كان
النسق \Log{G3c}~!!{prove}s نتيجة قاعدة منطقية، سوى
\RightR{\lexists} و\LeftR{\lforall}، فهو~!!{prove}s كل مقدمة من مقدمات
تلك القاعدة من غير أن يزيد !!{height} البرهان. بل يصح أن نقوي النتيجة،
فنثبت هذا الحكم نفسه في~$\Log{G3c} + \CutCS$.

ونجعل، كما تقدم، \emph{رتبة القطع} في استدلال \CutCS{} هي عمق
!!{formula} القطع الذي فيه.

\begin{lem}\ollabel{lem:inv-G3c-cut}
إذا كان $\Log{G3c} + \CutCS$~!!{prove}s نتيجة قاعدة منطقية~$\OLId{latin-looped}{R}$، ليست
\RightR{\lexists} ولا \LeftR{\lforall}، بـ!!a{proof}~$\pi$، فإنه
!!{prove}s كل مقدمة أيضًا بـ!!{proof}~$\pi'$ (أو بـ!!{proof}s~$\pi'$،
$\pi''$ إن كانت القاعدة ذات مقدمة واحدة أو مقدمتين). ومع ذلك يتحقق أن
(أ)~$\pheight{\pi'}\le\pheight{\pi}$ (وكذلك
$\pheight{\pi''}\le\pheight{\pi}$)، و(ب)~يكون عدد استدلالات \CutCS{}
ورتبها في~$\pi'$ (وفي $\pi''$) هو نفسه في~$\pi$.
\end{lem}

\begin{proof}
وبيان ذلك أن ننظر في براهين
\cref{pt:seq:inv:lem:G3c-invert,pt:seq:inv:lem:invert-quant}. فالاستقراء
فيها جار على~$\pheight{\pi}$. أما الأساس فيستبدل مسلّمة بمسلّمة، فيصح
الشرطان (أ) و(ب) بلا مزيد بيان. وإذا كانت آخر قاعدة في~$\pi$ هي القاعدة
المقصودة~$\OLId{latin-looped}{R}$، كانت !!{proof}sها الفرعية المباشرة، وهي
!!{proof}s~$\pi_\OLId{latin-ordinary}{i}$، عين المطلوب، وصح الشرطان كذلك. وأما سائر الحالات
فنعمل فرض الاستقراء في !!{proof}s~$\pi$ الفرعية المباشرة، أي المنتهية
بمقدمة أو مقدمات آخر استدلال~$\OLId{latin-looped}{S}$، ثم نعيد القاعدة نفسها~$\OLId{latin-looped}{S}$ على ما نتج.
ولما كان الشرطان (أ) و(ب) صحيحين في !!{proof}s البرهان الجديد الفرعية
المباشرة، صحا في البرهان كله. فلم يبق إلا أن يكون آخر الاستدلال~\CutCS.
ونبينه، كما فعلنا من قبل، بمثال واحد: لنفرض أن
$\pi$ هو !!a{proof} لنتيجة \LeftR{\land}، وهي $!B \land !C, \OLId{greek-variable}{Gamma}
\Sequent \OLId{greek-variable}{Delta}$، وأنه ينتهي بــ~\CutCS:
\[
\AxiomC{}
\RightLabel{$\pi_\OLMathNumeral{1}$}
\Deduce$!B \land !C, \OLId{greek-variable}{Gamma} \fCenter \OLId{greek-variable}{Delta}, !A$
\AxiomC{}
\RightLabel{$\pi_\OLMathNumeral{2}$}
\Deduce$!A, !B \land !C, \OLId{greek-variable}{Gamma} \fCenter \OLId{greek-variable}{Delta}$
\RightLabel{\CutCS}
\BinaryInf$!B \land !C, \OLId{greek-variable}{Gamma} \fCenter \OLId{greek-variable}{Delta}$
\DisplayProof
\]
ينطبق فرض الاستقراء على $\pi_\OLMathNumeral{1}$ و~$\pi_\OLMathNumeral{2}$ (وهما كذلك !!{proof}s لأمثلة
من نتيجة $\LeftR{\land}$)، ويعطينا !!{proof}s~$\pi_\OLMathNumeral{1}'$ و~$\pi_\OLMathNumeral{2}'$ لكل من
$!B, !C, \OLId{greek-variable}{Gamma} \fCenter \OLId{greek-variable}{Delta}, !A$ و$!A, !B, !C, \OLId{greek-variable}{Gamma} \fCenter
\OLId{greek-variable}{Delta}$ على الترتيب. ونجمعهما باستعمال \CutCS{} لنحصل على~$\pi'$:
\[
\AxiomC{}
\RightLabel{$\pi_\OLMathNumeral{1}'$}
\Deduce$!B, !C, \OLId{greek-variable}{Gamma} \fCenter \OLId{greek-variable}{Delta}, !A$
\AxiomC{}
\RightLabel{$\pi_\OLMathNumeral{2}'$}
\Deduce$!A, !B, !C, \OLId{greek-variable}{Gamma} \fCenter \OLId{greek-variable}{Delta}$
\RightLabel{\CutCS}
\BinaryInf$!B, !C, \OLId{greek-variable}{Gamma} \fCenter \OLId{greek-variable}{Delta}$
\DisplayProof
\]
وعلى هذا يظهر تحقق الشرطين (أ) و(ب).
\end{proof}

وهذا الوجه لا يستقيم إذا استعملنا \Cut{} مكان \CutCS؛ إذ تشتمل نتيجة
!!{proof} الأخير عندئذ على نسختين من~$!B, !C, \OLId{greek-variable}{Gamma}$ في المقدم ونسختين
من~$\OLId{greek-variable}{Delta}$ في التالي. وحينئذ لا بد من مقبولية الانكماش، مع أن برهان
\olref[seq][inv]{prop:G3c-cont-adm} نفسه يستعمل مبرهنة قابلية العكس، فيدور
الاستدلال.

ومن \cref{pt:cut:inv:lem:inv-G3c-cut} يستخرج برهان آخر لحذف القطع. ولسنا
فيه نزيل المواضع العليا لاستدلالات \CutCS{} واحدًا بعد واحد، بل نقصد
مواضع استدلالات \CutCS{} التي بلغت الرتبة القصوى. فنبدأ بـ!!a{proof}~$\pi$
في~$\Log{G3c} + \CutCS$، ونزيل منه قطعًا حيث وقع، وربما أقمنا مكانه قطوعًا
أدنى رتبة، ثم نكرر العمل حتى لا يبقى قطع. وإنما يشترط ألا يكون فوق القطع
المعالج قطع مساو له في الرتبة أو أعلى منه.

\begin{lem}\ollabel{lem:max-cut-red-G3c} ليكن~$\pi$~!!a{proof} في
$\Log{G3c}+\CutCS$ آخره استدلال~\CutCS{} رتبته~$\OLId{latin-ordinary}{n}$، ولا يستعمل فيما سواه
إلا قواعد~$\Log{G3c}$ واستدلالات \CutCS{} رتبتها~$<\OLId{latin-ordinary}{n}$. فحينئذ يوجد برهان
$\pi'$ للمتتالية النهائية عينها في~$\Log{G3c} + \CutCS$، وكل استدلال
\CutCS{} فيه رتبته~$<\OLId{latin-ordinary}{n}$.
\end{lem}

\begin{proof}
آخر !!{proof}~$\pi$ على هذا الرسم:
\[
\AxiomC{}
\RightLabel{$\pi_\OLMathNumeral{1}$}
\Deduce$\OLId{greek-variable}{Gamma} \fCenter \OLId{greek-variable}{Delta}, !A$
\AxiomC{}
\RightLabel{$\pi_\OLMathNumeral{2}$}
\Deduce$!A, \OLId{greek-variable}{Gamma} \fCenter \OLId{greek-variable}{Delta}$
\RightLabel{\CutCS}
\BinaryInf$\OLId{greek-variable}{Gamma} \fCenter \OLId{greek-variable}{Delta}$
\DisplayProof
\]
نميّز حالات بحسب صورة~$!A$.

لنفرض أن $!A$ ذرية. لاحظ أنه، بمقتضى الشرط القائل إن كل استدلال \CutCS{}
في~$\pi_\OLMathNumeral{1}$ و$\pi_\OLMathNumeral{2}$ لا بد أن تكون رتبته $< \depth{!A} = \OLMathNumeral{0}$، لا يجوز أن يوجد أي
استدلال \CutCS{} في~$\pi_\OLMathNumeral{1}$ أو~$\pi_\OLMathNumeral{2}$. نبيّن بالاستقراء على
$\pheight{\pi_\OLMathNumeral{2}}$ أن لـ$\OLId{greek-variable}{Gamma} \Sequent \OLId{greek-variable}{Delta}$ برهانًا خاليًا من القطع.
إذا كان $\pheight{\pi_\OLMathNumeral{2}}=\OLMathNumeral{0}$، فإن $!A, \OLId{greek-variable}{Gamma} \Sequent \OLId{greek-variable}{Delta}$ مسلّمة.
فإذا لم تكن $!A$ رئيسة، كان بعض $!C$ الذري !!a{element} في كل من~$\OLId{greek-variable}{Gamma}$
و$\OLId{greek-variable}{Delta}$، وكانت $\OLId{greek-variable}{Gamma} \Sequent \OLId{greek-variable}{Delta}$ نفسها مسلّمة. وإذا كانت $!A$
رئيسة، فإن $\OLId{greek-variable}{Delta}=\OLId{greek-variable}{Delta}', !A$. وفي هذه الحالة ينتهي $\pi_\OLMathNumeral{1}$ بــ
$\OLId{greek-variable}{Gamma} \Sequent \OLId{greek-variable}{Delta}', !A, !A$. ونحصل على برهان خال من القطع لــ
$\OLId{greek-variable}{Gamma} \Sequent \OLId{greek-variable}{Delta}', !A$ بتطبيق
\olref[seq][inv]{prop:G3c-cont-adm} (مقبولية الانكماش). وإذا كان
$\pheight{\pi_\OLMathNumeral{2}}>\OLMathNumeral{0}$، فإنه ينتهي بقاعدة~$\OLId{latin-looped}{S}$. خذ مقدمة من مقدمات تلك القاعدة؛
صورتها $!A, \OLId{greek-variable}{Gamma}', \Pi \Sequent \OLId{greek-variable}{Delta}', \Lambda$، حيث $\Pi$ و$\Lambda$
هما !!{formula}s الفعالة للقاعدة~$\OLId{latin-looped}{S}$. يعطي تطبيق
\olref[seq][adm]{prop:weak-G3c-adm} على~$\pi_\OLMathNumeral{1}$ و$\pi_\OLMathNumeral{2}$ !!{proof}s لكل من
$\OLId{greek-variable}{Gamma}, \OLId{greek-variable}{Gamma}', \Pi \Sequent \OLId{greek-variable}{Delta}, \OLId{greek-variable}{Delta}', \Lambda, !A$ و
$!A, \OLId{greek-variable}{Gamma}, \OLId{greek-variable}{Gamma}', \Pi \Sequent \OLId{greek-variable}{Delta}, \OLId{greek-variable}{Delta}', \Lambda$، وينطبق
عليهما فرض الاستقراء. وعلى هذا ينتج لنا !!a{proof} لــ$\OLId{greek-variable}{Gamma}, \OLId{greek-variable}{Gamma}', \Pi
\Sequent \OLId{greek-variable}{Delta}, \OLId{greek-variable}{Delta}', \Lambda$، لكل مقدمة من مقدمات~$\OLId{latin-looped}{S}$. وبتطبيق $\OLId{latin-looped}{S}$
ينتج لنا $\OLId{greek-variable}{Gamma}, \OLId{greek-variable}{Gamma} \Sequent \OLId{greek-variable}{Delta}, \OLId{greek-variable}{Delta}$، ثم يزوّدنا
\olref[seq][inv]{prop:G3c-cont-adm} بالبرهان الخالي من القطع لــ$\OLId{greek-variable}{Gamma}
\Sequent \OLId{greek-variable}{Delta}$.

إذا لم تكن $!A$ ذرية، نميّز حالات بحسب !!{main operator} لــ$!A$. وفي كل
حالة نطبّق \olref{lem:inv-G3c-cut} لنحصل على !!{proof}s لمقدمات قاعدتي
اليسار واليمين اللتين تكون $!A$ صيغتهما الرئيسة. ثم نتابع كما في الحالة~E
من برهان \olref[top]{lem:cut-adm-G3c}.
\begin{enumerate}
  \item $!A \ident \lnot !B$. ينتهي !!{proof}s~$\pi_\OLMathNumeral{1}$ و$\pi_\OLMathNumeral{2}$ بكل من
  $\OLId{greek-variable}{Gamma} \Sequent \OLId{greek-variable}{Delta}, \lnot !B$ و$\lnot !B, \OLId{greek-variable}{Gamma} \Sequent
  \OLId{greek-variable}{Delta}$. وبحسب \olref{lem:inv-G3c-cut} يوجد !!{proof}s~$\pi_\OLMathNumeral{1}'$ و
  $\pi_\OLMathNumeral{2}'$ لــ$!B, \OLId{greek-variable}{Gamma} \Sequent \OLId{greek-variable}{Delta}$ و$\OLId{greek-variable}{Gamma} \Sequent \OLId{greek-variable}{Delta},
  !B$. ويكون !!{proof}~$\pi'$ كما يأتي:
  \[
  \AxiomC{}
  \RightLabel{$\pi_\OLMathNumeral{2}'$}
  \Deduce$\OLId{greek-variable}{Gamma} \fCenter \OLId{greek-variable}{Delta}, !B$
  \AxiomC{}
  \RightLabel{$\pi_\OLMathNumeral{1}'$}
  \Deduce$!B, \OLId{greek-variable}{Gamma} \fCenter \OLId{greek-variable}{Delta}$
  \RightLabel{\CutCS}
  \BinaryInf$\OLId{greek-variable}{Gamma} \fCenter \OLId{greek-variable}{Delta}$
  \DisplayProof
  \]
  \item $!A \ident (!B \lor !C)$. ينتهي !!{proof}s~$\pi_\OLMathNumeral{1}$ و$\pi_\OLMathNumeral{2}$
  بكل من $\OLId{greek-variable}{Gamma} \Sequent \OLId{greek-variable}{Delta}, !B \lor !C$ و$!B \lor !C, \OLId{greek-variable}{Gamma}
  \Sequent \OLId{greek-variable}{Delta}$. وبحسب \olref{lem:inv-G3c-cut} (قابلية عكس
  \RightR{\lor} مطبقة على~$\pi_\OLMathNumeral{1}$)، يوجد !!a{proof}~$\pi_\OLMathNumeral{1}'$ لــ
  $\OLId{greek-variable}{Gamma} \Sequent \OLId{greek-variable}{Delta}, !B, !C$. وبحسب \olref{lem:inv-G3c-cut}
  (قابلية عكس \LeftR{\lor} مطبقة على~$\pi_\OLMathNumeral{2}$)، يوجد !!a{proof}~$\pi_\OLMathNumeral{2}'$
  لــ$!B, \OLId{greek-variable}{Gamma} \Sequent \OLId{greek-variable}{Delta}$ و!!a{proof}~$\pi_\OLMathNumeral{2}''$ لــ$!C, \OLId{greek-variable}{Gamma}
  \Sequent \OLId{greek-variable}{Delta}$. وبتطبيق \olref[seq][adm]{prop:weak-G3c-adm} على
  $\pi_\OLMathNumeral{2}''$ نحصل كذلك على !!a{proof}~$\pi_\OLMathNumeral{2}'''$ لــ$!C, \OLId{greek-variable}{Gamma} \Sequent
  \OLId{greek-variable}{Delta}, !B$. ويكون !!{proof}~$\pi'$:
  \[
  \AxiomC{}
  \RightLabel{$\pi_\OLMathNumeral{1}'$}
  \Deduce$\OLId{greek-variable}{Gamma} \fCenter \OLId{greek-variable}{Delta}, !B, !C$
  \AxiomC{}
  \RightLabel{$\pi_\OLMathNumeral{2}'''$}
  \Deduce$!C, \OLId{greek-variable}{Gamma} \fCenter \OLId{greek-variable}{Delta}, !B$
  \RightLabel{\CutCS}
  \BinaryInf$\OLId{greek-variable}{Gamma} \fCenter \OLId{greek-variable}{Delta}, !B$
  \AxiomC{}
  \RightLabel{$\pi_\OLMathNumeral{2}'$}
  \Deduce$!B, \OLId{greek-variable}{Gamma} \fCenter \OLId{greek-variable}{Delta}$
  \RightLabel{\CutCS}
  \BinaryInf$\OLId{greek-variable}{Gamma} \fCenter \OLId{greek-variable}{Delta}$
  \DisplayProof
  \]
  \item $!A \ident (!B \land !C)$. تمرين.
  \item $!A \ident (!B \lif !C)$. ينتهي !!{proof}s~$\pi_\OLMathNumeral{1}$ و$\pi_\OLMathNumeral{2}$
  بكل من $\OLId{greek-variable}{Gamma} \Sequent \OLId{greek-variable}{Delta}, !B \lif !C$ و$!B \lif !C, \OLId{greek-variable}{Gamma}
  \Sequent \OLId{greek-variable}{Delta}$. وبحسب \olref{lem:inv-G3c-cut} (قابلية عكس
  \RightR{\lif} مطبقة على~$\pi_\OLMathNumeral{1}$)، يوجد !!a{proof}~$\pi_\OLMathNumeral{1}'$ لــ
  $!B, \OLId{greek-variable}{Gamma} \Sequent \OLId{greek-variable}{Delta}, !C$. وبحسب \olref{lem:inv-G3c-cut}
  (قابلية عكس \LeftR{\lif} مطبقة على~$\pi_\OLMathNumeral{2}$)، يوجد !!a{proof}~$\pi_\OLMathNumeral{2}'$
  لــ$\OLId{greek-variable}{Gamma} \Sequent \OLId{greek-variable}{Delta}, !B$ و!!a{proof}~$\pi_\OLMathNumeral{2}''$ لــ$!C, \OLId{greek-variable}{Gamma}
  \Sequent \OLId{greek-variable}{Delta}$. وبتطبيق \olref[seq][adm]{prop:weak-G3c-adm} على
  $\pi_\OLMathNumeral{2}''$ نحصل كذلك على !!a{proof}~$\pi_\OLMathNumeral{2}'''$ لــ$!C, !B, \OLId{greek-variable}{Gamma}
  \Sequent \OLId{greek-variable}{Delta}$. ويكون !!{proof}~$\pi'$:
  \[
  \AxiomC{}
  \RightLabel{$\pi_\OLMathNumeral{2}'$}
  \Deduce$\OLId{greek-variable}{Gamma} \fCenter \OLId{greek-variable}{Delta}, !B$
  \AxiomC{}
  \RightLabel{$\pi_\OLMathNumeral{1}'$}
  \Deduce$!B, \OLId{greek-variable}{Gamma} \fCenter \OLId{greek-variable}{Delta}, !C$
  \AxiomC{}
  \RightLabel{$\pi_\OLMathNumeral{2}'''$}
  \Deduce$!C, !B, \OLId{greek-variable}{Gamma} \fCenter \OLId{greek-variable}{Delta}$
  \RightLabel{\CutCS}
  \BinaryInf$!B, \OLId{greek-variable}{Gamma} \fCenter \OLId{greek-variable}{Delta}$
  \RightLabel{\CutCS}
  \BinaryInf$\OLId{greek-variable}{Gamma} \fCenter \OLId{greek-variable}{Delta}$
  \DisplayProof
  \]
  \item $!A \ident \lforall[\OLId{latin-ordinary}{x}][!B(\OLId{latin-ordinary}{x})]$. ينتهي !!{proof}s~$\pi_\OLMathNumeral{1}$ و
  $\pi_\OLMathNumeral{2}$ بكل من $\OLId{greek-variable}{Gamma} \Sequent \OLId{greek-variable}{Delta}, \lforall[\OLId{latin-ordinary}{x}][!B(\OLId{latin-ordinary}{x})]$ و
  $\lforall[\OLId{latin-ordinary}{x}][!B(\OLId{latin-ordinary}{x})], \OLId{greek-variable}{Gamma} \Sequent \OLId{greek-variable}{Delta}$. وبحسب
  \olref{lem:inv-G3c-cut} يوجد !!a{proof}~$\pi_\OLMathNumeral{1}'$ لــ$\OLId{greek-variable}{Gamma} \Sequent
  \OLId{greek-variable}{Delta}, !B(\OLId{latin-ordinary}{c})$.

  تأمل الآن !!{proof}~$\pi_\OLMathNumeral{2}$. إنه ينتهي بــ$\lforall[\OLId{latin-ordinary}{x}][!B(\OLId{latin-ordinary}{x})],
  \OLId{greek-variable}{Gamma} \Sequent \OLId{greek-variable}{Delta}$. وبينما نتحرك صعودًا ابتداء من المتتالية النهائية
  لــ$\pi_\OLMathNumeral{2}$، نعلّم كل ورود لــ$\lforall[\OLId{latin-ordinary}{x}][!B(\OLId{latin-ordinary}{x})]$ على يسار متتالية «يقود
  إلى» ورود $\lforall[\OLId{latin-ordinary}{x}][!B(\OLId{latin-ordinary}{x})]$ في المتتالية النهائية لــ$\pi_\OLMathNumeral{2}$، أي:
  (أ)~علّم $\lforall[\OLId{latin-ordinary}{x}][!B(\OLId{latin-ordinary}{x})]$ في المتتالية النهائية لــ$\pi_\OLMathNumeral{2}$.
  (ب)~إذا وردت $\lforall[\OLId{latin-ordinary}{x}][!B(\OLId{latin-ordinary}{x})]$ في سياق نتيجة قاعدة وكانت معلّمة،
  فعلّم الورود المناظر في مقدمة القاعدة أو مقدماتها. (ج)~إذا كانت
  $\lforall[\OLId{latin-ordinary}{x}][!B(\OLId{latin-ordinary}{x})]$ رئيسة في نتيجة قاعدة \LeftR{\lforall} وكانت معلّمة،
  فعلّم الورودين الفعالين المناظرين لــ$\lforall[\OLId{latin-ordinary}{x}][!B(\OLId{latin-ordinary}{x})]$ و$!B(\OLId{latin-ordinary}{t})$ في
  المقدمة.

  تأمل الآن كل هذه الورودات المعلّمة لــ$\lforall[\OLId{latin-ordinary}{x}][!B(\OLId{latin-ordinary}{x})]$. لا يكون أي
  منها فعالًا في قاعدة غير \LeftR{\lforall} (فلا تُعلّم إلا !!{formula}s
  الفعالة لقاعدة \LeftR{\lforall}). احذف من~$\pi_\OLMathNumeral{2}$ كل الورودات المعلّمة
  لــ$\lforall[\OLId{latin-ordinary}{x}][!B(\OLId{latin-ordinary}{x})]$.
  إذا كان الناتج !!{proof} صحيحًا، فلم تكن الورودات المعلّمة لــ
  $\lforall[\OLId{latin-ordinary}{x}][!B(\OLId{latin-ordinary}{x})]$ فعالة قط، بل أتت كلها مباشرة من مسلّمات. وينتهي
  !!{proof} بــ$\OLId{greek-variable}{Gamma} \Sequent \OLId{greek-variable}{Delta}$، ويمكننا أن نستبدل به $\pi$.
  وبذلك نكون قد أزلنا قطعًا ذا رتبة قصوى.

  وإلا، فالناتج ليس برهانًا صحيحًا، لأننا حذفنا !!{formula}s
  $\lforall[\OLId{latin-ordinary}{x}][!B(\OLId{latin-ordinary}{x})]$ كانت فعالة في بعض استدلالات \LeftR{\lforall}.
  وقد حوّلنا هذه إلى «استدلالات» من الصورة
  \[
  \AxiomC{}
  \RightLabel{$\pi_\OLId{latin-ordinary}{t}$}
  \Deduce$!B(\OLId{latin-ordinary}{t}), \Pi \fCenter \Lambda$
  \RightLabel{\LeftR{\lforall}}
  \UnaryInf$\Pi \fCenter \Lambda$
  \DisplayProof
  \]
  استبدل بكل «استدلال» علوي من هذا النوع ما يأتي:
  \[
  \AxiomC{}
  \RightLabel{$\Subst{\pi_\OLMathNumeral{1}'}{\OLId{latin-ordinary}{t}}{\OLId{latin-ordinary}{c}}[\Pi \Sequent \Lambda]$}
  \Deduce$\OLId{greek-variable}{Gamma}, \Pi \fCenter \OLId{greek-variable}{Delta}, \Lambda, !B(\OLId{latin-ordinary}{t})$
  \AxiomC{}
  \RightLabel{$\pi_\OLId{latin-ordinary}{t}[\OLId{greek-variable}{Gamma} \Sequent \OLId{greek-variable}{Delta}]$}
  \Deduce$!B(\OLId{latin-ordinary}{t}), \OLId{greek-variable}{Gamma}, \Pi \fCenter \OLId{greek-variable}{Delta}, \Lambda$
  \RightLabel{\CutCS}
  \BinaryInf$\OLId{greek-variable}{Gamma}, \Pi \fCenter \OLId{greek-variable}{Delta}, \Lambda$
  \DisplayProof
  \]
  وأضف $\OLId{greek-variable}{Gamma}$ إلى يسار كل متتالية أسفله، و$\OLId{greek-variable}{Delta}$ إلى يمينها. كرر ذلك
  حتى تُستبدل جميع «استدلالات» \LeftR{\lforall} التي تحتوي مقدمتها على
  $!B(\OLId{latin-ordinary}{t})$ معلّمة باستدلال \CutCS{} على $!B(\OLId{latin-ordinary}{t})$ ما. وتكون المتتالية النهائية
  لــ!!{proof} الناتج (وهو الآن !!{proof} صحيح) من الصورة
  $\OLId{greek-variable}{Gamma}, \dots, \OLId{greek-variable}{Gamma} \Sequent \OLId{greek-variable}{Delta}, \dots, \OLId{greek-variable}{Delta}$. طبّق مقبولية
  الانكماش لتحوّله إلى !!a{proof} لــ$\OLId{greek-variable}{Gamma} \Sequent \OLId{greek-variable}{Delta}$، واستبدل به
  $\pi$. وعلى هذا استبدلنا قطعًا واحدًا على $\lforall[\OLId{latin-ordinary}{x}][!B(\OLId{latin-ordinary}{x})]$ بقطوع
  (قد تكون كثيرة) على $!B(\OLId{latin-ordinary}{t})$، لكنها كلها أدنى رتبة. وتبقى أي قطوع كانت
  موجودة أصلًا في $\pi_\OLMathNumeral{1}$ أو $\pi_\OLMathNumeral{2}$، لكنها كذلك كلها أدنى رتبة.
\end{enumerate}
\end{proof}

\begin{prob}
تحقق من الحالة $!A \ident (!B \land !C)$ في برهان
\olref{lem:max-cut-red-G3c}. أي بيّن أنه إذا وجد !!a{proof}~$\pi$ ينتهي بــ
\[
\AxiomC{}
\RightLabel{$\pi_\OLMathNumeral{1}$}
\Deduce$\OLId{greek-variable}{Gamma} \fCenter \OLId{greek-variable}{Delta}, !B \land !C$
\AxiomC{}
\RightLabel{$\pi_\OLMathNumeral{2}$}
\Deduce$!B \land !C, \OLId{greek-variable}{Gamma} \fCenter \OLId{greek-variable}{Delta}$
\RightLabel{\CutCS}
\BinaryInf$\OLId{greek-variable}{Gamma} \fCenter \OLId{greek-variable}{Delta}$
\DisplayProof
\]
وكان كل من $\pi_\OLMathNumeral{1}$ و$\pi_\OLMathNumeral{2}$ لا يحتوي إلا قطوعًا رتبتها
$<\depth{!B \land !C}$، فهناك !!a{proof}~$\pi'$ لــ$\OLId{greek-variable}{Gamma} \Sequent
\OLId{greek-variable}{Delta}$ لا يحتوي إلا قطوعًا رتبتها $<\depth{!B \land !C}$.
\end{prob}

% \begin{prob}
% تحقق من الحالة $!A \ident \lexists[x][!B(x)]$ في برهان
% \olref{lem:inv-G3c-cut}.
% \end{prob}


تثبت \olref{lem:max-cut-red-G3c} أنه يمكننا استبدال استدلال \CutCS{} ذي
الرتبة الأعلى في !!a{proof} باستدلالات \CutCS{} أدنى رتبة. ويمكننا مرة
أخرى استعمال هذه المبرهنة لإزالة أي عدد من استدلالات \CutCS{} من
!!a{proof}، وذلك بتطبيقها تباعًا على !!{proof}s الفرعية العليا المنتهية
بــ\CutCS{} ذي الرتبة القصوى.

\begin{thm}\ollabel{thm:cut-elim-G3c-largest}
يمكن تحويل أي !!{proof}~$\pi$ في $\Log{G3c} + \CutCS$ إلى برهان في~\Log{G3c}.
\end{thm}

\begin{proof}
  نثبت أولًا أنه يمكن إزالة كل القطوع ذات الرتبة القصوى في~$\pi$. لتكن
  $\OLId{latin-ordinary}{d}$ الرتبة القصوى للقطوع الواقعة في~$\pi$، وليكن $\OLId{latin-ordinary}{n}$ عدد القطوع ذات
  الرتبة~$\OLId{latin-ordinary}{d}$. البرهان بالاستقراء على~$\OLId{latin-ordinary}{n}$. إذا كان $\OLId{latin-ordinary}{n}=\OLMathNumeral{0}$ فلا شيء نثبته.
  وإذا كان $\OLId{latin-ordinary}{n}>\OLMathNumeral{0}$، فاختر قطعًا علويًا رتبته~$\OLId{latin-ordinary}{d}$، وليكن $\pi_\OLMathNumeral{1}$ هو
  !!{proof} الفرعي المنتهي به. وبحسب \olref{lem:max-cut-red-G3c} يوجد
  !!a{proof}~$\pi_\OLMathNumeral{1}'$ للمتتالية النهائية نفسها، وكل قطوعه رتبتها $<\OLId{latin-ordinary}{d}$.
  استبدل $\pi_\OLMathNumeral{1}$ بــ$\pi_\OLMathNumeral{1}'$ في~$\pi$. ينتج من ذلك !!a{proof} يحتوي
  $\OLId{latin-ordinary}{n}-\OLMathNumeral{1}$ قطعًا من الرتبة~$\OLId{latin-ordinary}{d}$، فينطبق فرض الاستقراء.

  والآن تتبع النتيجة بالاستقراء على~$\OLId{latin-ordinary}{d}$. فإذا كان $\OLId{latin-ordinary}{d}=\OLMathNumeral{0}$، كانت جميع القطوع
  ذرية، ويمكن إزالة جميعها بحسب النتيجة السابقة. وإذا كان $\OLId{latin-ordinary}{d}>\OLMathNumeral{0}$، تعطينا
  النتيجة السابقة برهانًا استُبدلت فيه كل القطوع ذات الرتبة~$\OLId{latin-ordinary}{d}$ بقطوع
  رتبتها~$<\OLId{latin-ordinary}{d}$. ولما كانت $\OLId{latin-ordinary}{d}$ هي الرتبة القصوى للقطوع في~$\pi$، أصبح لدينا
  برهان رتبة قطوعه القصوى $<\OLId{latin-ordinary}{d}$، فينطبق فرض الاستقراء.
\end{proof}

\end{document}
